\section{Introduction}
\label{sec:intro}

% Para 1: Problem — VLA success + robustness gap under distribution shift
Vision-language-action models (VLAs) that generate actions via flow matching~\cite{black2024pi0,shukor2025smolvla} have reached 97--99\% success rates on standard manipulation benchmarks such as LIBERO~\cite{liu2023libero}.
Yet this near-saturation masks a fragile reality: under distribution shifts common in real-world deployment---camera viewpoint changes, lighting variation, sensor noise, object layout differences---the same policies suffer substantial performance degradation~\cite{fei2025liberoplus}.
Closing this robustness gap is a prerequisite for deploying learned manipulation policies outside tightly controlled lab settings.

% Para 2: Existing approaches — training-time vs inference-time, what they miss
Two lines of work address VLA robustness.
\emph{Training-time methods} improve the base policy itself: SRPO~\cite{srpo2025} applies self-referential policy optimization, while RobustVLA~\cite{robustvla2025} and DiG-Flow~\cite{digflow2025} incorporate robustness-aware objectives into training.
These methods are effective but require retraining or fine-tuning, and their gains are bounded by the training distribution.
\emph{Inference-time methods} leave the policy frozen and intervene at deployment.
ReconVLA~\cite{chen2026reconvla} applies conformal quantile regression to calibrate action-level uncertainty bounds and detects out-of-distribution states via squared Mahalanobis distance---an uncertainty quantification framework, not an action improvement method.
SAFE~\cite{gu2025safe} detects anomalous states and triggers fallback behaviors.
UPS~\cite{yuan2026ups} and CP-in-the-Loop~\cite{wang2025cploop} use conformal prediction for binary decisions---execute or abstain, safe or unsafe---but do not use conformal statistics as features for continuous correction.
A2C2~\cite{a2c2_2025} learns a residual correction that runs asynchronously alongside the base policy, yet applies it without uncertainty conditioning, lacking any mechanism to modulate correction magnitude across action dimensions.
In short, no existing method performs inference-time action correction conditioned on \emph{structured, per-degree-of-freedom uncertainty}.

% Para 3: Our insight — per-DOF conformal features as continuous input to learned correction
Our key observation is that the stochastic samples already produced by a flow-matching VLA contain rich, structured uncertainty information that current methods discard.
By fitting a shrinkage-regularized Gaussian to $K{=}10$ samples and computing multivariate conformal statistics~\cite{braun2025cgs}, we obtain a 40-dimensional feature vector that encodes \emph{which} action dimensions are uncertain (per-DOF standard deviations), the \emph{calibrated confidence level} (conformal radii from Adaptive Conformal Inference~\cite{gibbs2021aci}), and the \emph{overall predictive spread} (log-volume of the covariance ellipsoid).
This representation captures substantially more structure than the scalar distances used in ReconVLA~\cite{chen2026reconvla}: it preserves the cross-DOF covariance structure that enables a downstream module to correct high-uncertainty dimensions aggressively while leaving confident predictions untouched.
A controlled 3D point-reaching experiment confirms this: a correction MLP conditioned on per-DOF uncertainty reduces total error by 6.7\% over scalar conditioning and by 16.7\% over no uncertainty input, with 22--24\% improvements concentrated on the non-perturbed axes where scalar conditioning cannot distinguish signal from noise (Section~\ref{sec:motivating}).

% Para 4: ACR-VLA overview — 3 components, key design choices
We introduce \textbf{ACR-VLA} (Adaptive Conformal Refinement for VLA), an inference-time module that attaches to any frozen flow-matching VLA and performs dimension-selective action correction.
ACR-VLA comprises three components:
(i)~\textbf{Conformalized Gaussian Scoring (CGS)} draws $K{=}10$ samples per control step, fits a Ledoit--Wolf shrinkage covariance, and extracts a 40D structured uncertainty feature;
(ii)~\textbf{Adaptive Conformal Inference with Chunk Overlap Consistency (ACI+COC)} maintains an online conformal threshold~\cite{gibbs2021aci} using a self-referential coverage signal derived from overlapping action chunks---requiring no external supervisor, reward model, or ground-truth actions;
(iii)~\textbf{Conformal Refinement Module (CRM)}, a lightweight MLP (${\sim}$337K parameters) that takes the CGS features as conditioning and outputs a dimension-selective residual correction.
We note that conformal prediction serves here as a \emph{feature extraction and calibration mechanism}, not as a guarantee on the corrected output: CGS features provide structured uncertainty conditioning for CRM, while ACI+COC adapts these features online under distribution shift.
We instantiate ACR-VLA on SmolVLA~\cite{shukor2025smolvla}, a 450M-parameter flow-matching VLA.

% Para 5: Contributions — numbered list
Our contributions are:
\begin{enumerate}[leftmargin=*,nosep]
  \item \textbf{Per-DOF conformal features for dimension-selective correction.}
  We show that the 40D CGS representation---per-DOF standard deviations, calibrated conformal radii, and covariance log-volume---enables CRM to selectively correct uncertain action dimensions while preserving accurate ones.
  On LIBERO, ACR-VLA improves base SmolVLA success by \textbf{XX.X}\,pp, and ablations confirm that per-DOF features outperform both scalar uncertainty and no-uncertainty baselines (Section~\ref{sec:experiments}).

  \item \textbf{Self-referential online calibration via Chunk Overlap Consistency.}
  COC exploits the temporal overlap in receding-horizon action chunks to provide a coverage signal for ACI without external supervision.
  This enables the conformal threshold to adapt online under arbitrary distribution shift, maintaining calibrated uncertainty features throughout deployment (Section~\ref{sec:aci}).

  \item \textbf{Systematic evaluation across multiple LIBERO task suites.}
  We evaluate ACR-VLA on LIBERO-Spatial and LIBERO-Object, providing per-task breakdowns, ablations isolating the contribution of each component (CGS features, ACI+COC calibration, CRM architecture), and analysis of when and why dimension-selective correction helps (Section~\ref{sec:experiments}).
\end{enumerate}
